\documentclass[11pt,letterpaper]{article}
\usepackage[margin=1in]{geometry}
\usepackage[T1]{fontenc}
\usepackage{newtxtext,newtxmath}
\usepackage{microtype}
\usepackage[hidelinks]{hyperref}

\hypersetup{
  pdftitle={Conditioning odor-identity RSA: Methods},
  pdfsubject={Proposed Methods addition for PI review}
}
\setlength{\parindent}{0pt}
\setlength{\parskip}{0.8\baselineskip}
\pagestyle{empty}

\begin{document}

{\small\sffamily Draft for PI review \hfill November 12, 2026\par}
\vspace{0.4\baselineskip}
{\Large\bfseries Conditioning odor-identity RSA: Methods\par}
\vspace{0.4\baselineskip}

Conditioning single-trial GLMs used the same preprocessed images and nuisance regressors as the probe models. Each least-squares-separate model contained a target-cue regressor spanning the cue and prediction period (5\,s) and a second regressor pooling all other cue presentations with the same duration. Additional regressors modeled the delay (2\,s, beginning 5\,s after cue onset), odor delivery (2\,s, beginning 7\,s after cue onset), zero-duration button presses at the recorded response times, and fixation replacing delay and delivery on missed-response trials (4\,s). Delivery events for both food odors and clean air were pooled into one regressor and included on all responded trials, including incorrect predictions. All cue presentations were modeled; correct-trial filtering was applied subsequently for RSA.

Odor-identity RSA used the same bilateral OFC regions and procedures as the probe analysis. Only correct food-odor trials with recorded responses and finite positive response times were retained, pooling all available conditioning runs. Two additional participants were excluded for sparse correct trials and severe OFC coverage loss, leaving 29 participants; 28 contributed three runs and one contributed runs 1--2 because the third scan was interrupted. Participant-level RSA coefficients were assessed using one-sided one-sample $t$-tests against zero, with uncorrected $p$ values.

\end{document}
